% ============================================================
% Section 4: Numerical Verification and Results Interpretation
% ============================================================

\section{Numerical Verification and Results Interpretation}
\pagecheck{4--5}

\instructionplaceholder{Result strategy}{Present evidence in two stages: (i) numerical verification (convergence, stability, and modal checks), then (ii) interpreted engineering results under relevant operating conditions. Do not present final engineering conclusions before verification evidence is shown.}

\subsection{Numerical verification}
\instructionplaceholder{Required evidence}{Summarise verification outcomes for the implemented FEM model: consistency checks, benchmark comparisons against literature data, and error trends that justify the credibility of the final simulations. Validate structural deformations and displacements against published OC3/NREL 5MW reference results.}

\subsection{Convergence and time-step assessment}
\instructionplaceholder{Required evidence}{Provide mesh (spatial) and time-step sensitivity analysis. Quantify numerical uncertainty and justify the selected discretisation levels used for all reported final results.}

\subsection{Modal analysis}
\instructionplaceholder{Required evidence}{Identify natural frequencies and mode shapes of the assembled structural model. Relate dominant modes to expected forced-response characteristics under wave, wind, and current loading. Compare computed natural frequencies with literature values where available.}

\subsection{Dynamic response of the structure}
\instructionplaceholder{Required evidence}{Present the dynamic response at key locations (nacelle, centre of gravity, fairlead, etc.) under the operating conditions specified in the project description. Include time histories, frequency spectra, and maximum response estimates as appropriate. Discuss the physical interpretation of the results.}

\subsection{Reduced-order model response}
\instructionplaceholder{Required evidence}{Apply the reduced-order model using a selection of relevant mode shapes. Compare the reduced-order response against the full-model solution and discuss the accuracy and efficiency of the reduction.}

\subsection{Limitations and improvements}
\instructionplaceholder{Required evidence}{State model limitations and propose concrete improvements to geometry, loading fidelity, boundary conditions, or numerical schemes for future work.}
